% GENERATED FILE. Edit canonical lessons, then run npm run generate:books.
% canonical-source-hash: fnv1a64:efaff23d68806df4
% canonical-lessons: GU-C283-jo-to, GU-C283-jo-varsad-padshe, GU-C283-nahintar, GU-R283-if-and-otherwise

\chapter{If, then, or else}
\label{ch:gu-if-then-or-else}
\hlchaptermodality{283}

\begin{chapteropening}
\textbf{By the end of this chapter:} \emph{I can say what will happen if something does (jo varsād paḍśe, to …), give advice on a condition, and warn with nahī̃tar, or else.}

Complete GU-R283-if-and-otherwise: say what happens if, and what happens otherwise.
\end{chapteropening}

\section[\texorpdfstring{\gu{જો} \gu{તમને} \gu{તાવ} \gu{છે}, \gu{તો} … (jo tamne tāv …)}{જો તમને તાવ છે, તો … (jo tamne tāv …)}]{\gu{જો} …, \gu{તો} … (jo …, to …) — if, then}

\label{lesson:GU-C283-jo-to}



\begin{quote}
Say \textquotedblleft{}if\textquotedblright{}, then \textquotedblleft{}if … then\textquotedblright{}, from A2. (\textbf{\gu{જો}}, \textbf{\gu{જો} … \gu{તો} …}.) Then \textquotedblleft{}you should rest\textquotedblright{}. (\emph{tamāre ārām karvo joīe}.)
\end{quote}

\begin{grammarlens}[title={Grammar Lens: two halves}]
\begin{quote}
\textbf{\gu{જો} \gu{તમને} \gu{તાવ} \gu{છે}, \gu{તો} \gu{તમારે} \gu{આરામ} \gu{કરવો} \gu{જોઈએ}.}
\end{quote}

\begin{quote}
\emph{jo tamne tāv che, to tamāre ārām karvo joīe.} — \textquotedblleft{}If you have a fever, then you should rest.\textquotedblright{}
\end{quote}

\textbf{\gu{જો}} (\emph{jo}) opens the condition, and \textbf{\gu{તો}} (\emph{to}), \textquotedblleft{}then\textquotedblright{}, opens what follows from it. Each half is a whole sentence with its own verb. In quick speech \textbf{\gu{તો}} can drop, but written down it marks where the second half begins.
\end{grammarlens}

\subsection*{You'll want to know: if it is so now}
\begin{quote}
\textbf{\gu{જો} \gu{તમને} \gu{ભૂખ} \gu{છે}, \gu{તો} \gu{ખાઈ} \gu{લો}.}
\end{quote}

\begin{quote}
\emph{jo tamne bhūkh che, to khāī lo.} — \textquotedblleft{}If you're hungry, then go ahead and eat.\textquotedblright{}
\end{quote}

When the condition is about now, it takes the present: \textbf{\gu{જો} \gu{તમને} \gu{ભૂખ} \gu{છે}}, \textbf{\gu{જો} \gu{તે} \gu{ઘરે} \gu{છે}}. The second half can be advice, a request with \textbf{-\gu{ઓ}}, or a plain statement: \textbf{\gu{જો} \gu{ઠંડી} \gu{છે}, \gu{તો} \gu{હું} \gu{ઘરે} \gu{રહીશ}} (\emph{jo ṭhanḍī che, to hũ ghare rahīś}).

\subsection*{Your turn}
Say these aloud:

\begin{itemize}
\raggedright
  \item \textquotedblleft{}If you're hungry, go ahead and eat\textquotedblright{} — \emph{jo tamne bhūkh che, to khāī lo}
  \item \textquotedblleft{}If it's cold, I'll stay at home\textquotedblright{} — \emph{jo ṭhanḍī che, to hũ ghare rahīś}
\end{itemize}

\begin{tcolorbox}[breakable,colback=teal!4,colframe=teal!35!black,title={Before you move on}]
Which word opens the condition? (\textbf{\gu{જો}}.) Which opens the second half? (\textbf{\gu{તો}}.) \textquotedblleft{}If you have a fever\textquotedblright{}? (\textbf{\gu{જો} \gu{તમને} \gu{તાવ} \gu{છે}}.)
\end{tcolorbox}
\section[\texorpdfstring{\gu{જો} \gu{વરસાદ} \gu{પડશે}, \gu{તો} … (jo varsād paḍśe …)}{જો વરસાદ પડશે, તો … (jo varsād paḍśe …)}]{\gu{જો} \gu{વરસાદ} \gu{પડશે}, \gu{તો} … (jo varsād paḍśe, to …) — if it happens}

\label{lesson:GU-C283-jo-varsad-padshe}



\begin{quote}
Say \textquotedblleft{}if you have a fever, then you should rest\textquotedblright{}. (\emph{jo tamne tāv che, to tamāre ārām karvo joīe}.) Then \textquotedblleft{}we will stay\textquotedblright{}. (\emph{ame rahīśũ}.)
\end{quote}

\begin{grammarlens}[title={Grammar Lens: the future after \gu{જો}}]
\begin{quote}
\textbf{\gu{જો} \gu{વરસાદ} \gu{પડશે}, \gu{તો} \gu{અમે} \gu{ઘરે} \gu{રહીશું}.}
\end{quote}

\begin{quote}
\emph{jo varsād paḍśe, to ame ghare rahīśũ.} — \textquotedblleft{}If it rains, we'll stay at home.\textquotedblright{}
\end{quote}

English says \textquotedblleft{}if it rains\textquotedblright{}; Gujarati says \textquotedblleft{}if it will rain\textquotedblright{}: a condition about later takes the future, \textbf{\gu{પડશે}}, and so does the result, \textbf{\gu{રહીશું}}. \textquotedblleft{}If it doesn't rain\textquotedblright{} is the \textbf{\gu{નહીં}} form: \textbf{\gu{જો} \gu{વરસાદ} \gu{નહીં} \gu{પડે}, \gu{તો} \gu{અમે} \gu{જઈશું}} (\emph{jo varsād nahī̃ paḍe, to ame jaīśũ}).
\end{grammarlens}

\subsection*{You'll want to know: the past after \gu{જો}}
\begin{quote}
\textbf{\gu{જો} \gu{ટ્રેન} \gu{છૂટી} \gu{ગઈ}, \gu{તો} \gu{હું} \gu{બસમાં} \gu{આવીશ}.}
\end{quote}

\begin{quote}
\emph{jo ṭren chūṭī gaī, to hũ basmā̃ āvīś.} — \textquotedblleft{}If I miss the train, I'll come by bus.\textquotedblright{}
\end{quote}

When the condition would be over before the result starts, it can take the plain past: \textquotedblleft{}if the train has gone, then …\textquotedblright{}. \textbf{\gu{જો} \gu{મારું} \gu{પાકીટ} \gu{ખોવાઈ} \gu{ગયું}, \gu{તો} \gu{હું} \gu{પોલીસમાં} \gu{ફરિયાદ} \gu{કરીશ}} (\emph{jo mārũ pākīṭ khovāī gayũ, to hũ polīsmā̃ phariyād karīś}).

\subsection*{Your turn}
Say these aloud:

\begin{itemize}
\raggedright
  \item \textquotedblleft{}If it doesn't rain, we'll go\textquotedblright{} — \emph{jo varsād nahī̃ paḍe, to ame jaīśũ}
  \item \textquotedblleft{}If my wallet is lost, I'll make a complaint\textquotedblright{} — \emph{jo mārũ pākīṭ khovāī gayũ, to hũ phariyād karīś}
\end{itemize}

\begin{tcolorbox}[breakable,colback=teal!4,colframe=teal!35!black,title={Before you move on}]
\textquotedblleft{}If it rains\textquotedblright{}? (\textbf{\gu{જો} \gu{વરસાદ} \gu{પડશે}}.) \textquotedblleft{}If it doesn't rain\textquotedblright{}? (\textbf{\gu{જો} \gu{વરસાદ} \gu{નહીં} \gu{પડે}}.) \textquotedblleft{}If I miss the train\textquotedblright{}? (\textbf{\gu{જો} \gu{ટ્રેન} \gu{છૂટી} \gu{ગઈ}}.)
\end{tcolorbox}
\section[\texorpdfstring{\gu{નહીંતર} (nahī̃tar)}{નહીંતર (nahī̃tar)}]{\gu{નહીંતર} (nahī̃tar) — otherwise; or else}

\label{lesson:GU-C283-nahintar}



\begin{quote}
Before the new one: say \textquotedblleft{}if it rains, we'll stay at home\textquotedblright{}. (\emph{jo varsād paḍśe, to ame ghare rahīśũ}.) Then \textquotedblleft{}quickly\textquotedblright{} and \textquotedblleft{}an umbrella\textquotedblright{}. (\textbf{\gu{જલદી}}, \textbf{\gu{છત્રી}}.)
\end{quote}

\subsection*{You'll want to know: \gu{નહીંતર}}
\textbf{\gu{નહીંતર}} — \emph{nahī̃tar} — \textquotedblleft{}otherwise; or else\textquotedblright{}. First what to do, then \textbf{\gu{નહીંતર}}, then what will happen if it is not done:

\begin{quote}
\textbf{\gu{જલદી} \gu{ચાલો}, \gu{નહીંતર} \gu{ટ્રેન} \gu{છૂટી} \gu{જશે}.}
\end{quote}

\begin{quote}
\emph{jaldī cālo, nahī̃tar ṭren chūṭī jaśe.} — \textquotedblleft{}Walk quickly, or else we'll miss the train.\textquotedblright{}
\end{quote}

\begin{quote}
\textbf{\gu{છત્રી} \gu{લો}, \gu{નહીંતર} \gu{તમે} \gu{પલળી} \gu{જશો}.}
\end{quote}

\begin{quote}
\emph{chatrī lo, nahī̃tar tame palḷī jaśo.} — \textquotedblleft{}Take an umbrella, otherwise you'll get soaked.\textquotedblright{}
\end{quote}

It says in one word what \textbf{\gu{જો} … \gu{નહીં} …, \gu{તો} …} says in a whole clause.

\begin{cousinweb}
\textbf{\gu{નહીંતર}} (\emph{nahī̃tar}) is built from \textbf{\gu{નહીં}}, \textquotedblleft{}not\textquotedblright{}, and \emph{-tar}: \textquotedblleft{}if not, then\textquotedblright{}. CDIAL lists Gujarati \emph{nahitar}, \textquotedblleft{}otherwise\textquotedblright{}, under Sanskrit \emph{tarhi}, \textquotedblleft{}then\textquotedblright{} (CDIAL 5730).
\end{cousinweb}

\subsection*{Your turn}
Say these aloud:

\begin{itemize}
\raggedright
  \item \emph{nahī̃tar}
  \item \textquotedblleft{}Walk quickly, or else we'll miss the train\textquotedblright{} — \emph{jaldī cālo, nahī̃tar ṭren chūṭī jaśe}
  \item \textquotedblleft{}Take an umbrella, otherwise you'll get soaked\textquotedblright{} — \emph{chatrī lo, nahī̃tar tame palḷī jaśo}
\end{itemize}

\begin{tcolorbox}[breakable,colback=teal!4,colframe=teal!35!black,title={Before you move on}]
What does \textbf{\gu{નહીંતર}} mean? (\textquotedblleft{}Otherwise; or else\textquotedblright{}.) Which two words are inside it? (\textbf{\gu{નહીં}} and \emph{-tar}, \textquotedblleft{}then\textquotedblright{}.)
\end{tcolorbox}
\section[(review)]{Review — if, then, or else}

\label{lesson:GU-R283-if-and-otherwise}



\begin{quote}
Say \textquotedblleft{}the city is bigger than the village\textquotedblright{}. (\emph{śaher gām kartā̃ moṭũ che}.) Then \textquotedblleft{}I have lived here for two years\textquotedblright{}. (\emph{hũ be varasthī ahī̃ rahũ chũ}.) And \textquotedblleft{}my wish is to live in the village\textquotedblright{}. (\emph{mārī icchā gāmmā̃ rahevānī che}.)
\end{quote}

\subsection*{Your turn}
Say these aloud:

\begin{itemize}
\raggedright
  \item If you have a fever, then you should rest.
  \item If you're hungry, go ahead and eat.
  \item If it rains, we'll stay at home. If it doesn't rain, we'll go.
  \item If I miss the train, I'll come by bus. The bus is the cheapest.
  \item The shop is open from morning until evening. Go quickly, otherwise it will close.
\end{itemize}

\begin{tcolorbox}[breakable,colback=teal!4,colframe=teal!35!black,title={Before you move on}]
\textquotedblleft{}If it rains\textquotedblright{}, about tomorrow: which tense? (The future: \textbf{\gu{પડશે}}.) \textquotedblleft{}If I miss the train\textquotedblright{}? (\textbf{\gu{જો} \gu{ટ્રેન} \gu{છૂટી} \gu{ગઈ}}.) \textquotedblleft{}Or else\textquotedblright{}? (\textbf{\gu{નહીંતર}}.)
\end{tcolorbox}
